\documentclass[11pt]{article}
\usepackage{geometry}
\geometry{
letterpaper,
top=1in,
left=1in,
right=1in,
bottom=1in
}

% required packages
\usepackage[hidelinks]{hyperref}
\usepackage{titlesec}
\usepackage{color}
\usepackage{verbatim}
\usepackage{enumitem}
\usepackage{tabularx}

% line spacing
\usepackage{setspace}
\setlength{\parindent}{20pt}
\setlength{\parskip}{1em}

% margin adjustments
\raggedbottom

% title stuff
\title{How many words do you know in your native language?}
\author{Rey Stone}

% ====================
% DOCUMENT STARTS HERE
% ====================
\begin{document}

\onehalfspacing

\maketitle

\section {What is my "native language?"}

Before we get to the core question, I would like to define "native language". A quick google search of the query "what is someone's native language" gives me many search results, the first two being Wikipedia and Reddit. I love the internet. Wikipedia's definition of a native language is "A first language (L1), native language, native tongue, or mother tongue is the first language a person has been exposed to from birth or within the critical period". This raises some very interesting questions for me, because the primary language I speak is English, yet that was not the first language I was exposed to, nor was it the only one I learned growing up as a child. The first language I was exposed to was Sichuanese, a dialect of Mandarin, when I was living in an orphanage for the first thirteen months of my life. Afterwards, I was adopted by my parents and they taught me English for the first five years of my life. My family and I proceeded to move to China for job reasons, and I started learning Standard Mandarin (spoken and written) while attending international school. I moved back to the States when I was around 10 years old, thus solidifying my language of choice as English. In short, if you consider my "native language" the one I was exposed to first, then I know roughly 50 words in Mandarin and 0 in Sichuanese.

\section{How do I "know" a word?}

Already this assignment has confounded me in ways not anticipated. However, for simplicity's sake, and one that is probably the most accurate to the definition of "native language", my native language is English. Now that that's out of the way I'm going to get back to the core question. The first thing I have done is run this exact document using the methods shown in class. The output of that file was 159 words! So I know about 159 words confirmed.

But how do I "know" these words? Just because I can string together a coherent sentence doesn't necessarily mean that I "know" these words. Let's use some philosophy to define how I "know" a word. The branch of philosophy that deals with the theory of knowledge is epistemology which describes the nature of knowledge and the different types of knowledge. According to epistemology, knowledge is an awareness, understanding, familiarity or skill. To "have knowledge" one must make successful contact with reality. Through this contact, we can understand, interpret, and interact with the world. Of course, there are many counterarguments to what knowledge really is and how we can ever know anything, or if there is no knowledge at all. Thus, for this assignment, and before I go down too many rabbit holes of what it means to "know" something, knowing a word means that I can do one of the following:


\begin{itemize}
  \item Define it by recalling its dictionary definition
  \item Interact with it by using it in a sentence that others can understand and interpret themselves
\end{itemize}

I am intentionally leaving out writing, because while reading and writing are good markers of literacy, knowing words in a native language does not have to include literacy rates. There are many people who can go throughout their world without knowing how to spell half of the words they understand through experiencing the world (refer to any reddit post available on the internet.)

\section{What counts as an English "word"?}

Now that I have somewhat come to an arbitrary way of what it means to "know a word", we can focus on the question "what counts as an English word?". An agreed upon answer of what a word is a "minimum free form -- the smallest meaningful unit of speech that can function independently in a sentence". Therefore I will define a word as being: \textit{a lemma that function as a single semantic and grammatical unit, excluding derivations that would count as separate word families}.


\section{How many words do I know in English?}

Finally I can ascertain "how many?" by doing several online tests when I Google "how many words do you know in english test". Keeping in mind the criteria I defined for these tests above, since I am self-administering and testing on myself, I will try my best to adhere to the conditions defined above.

After about 5 tests that took a total of 20 minutes, a rough estimate that I consistently got was about 20,000 the lowest being 15,000 and the highest being 30,000. 2 out of 5 of these tests were essentially "pick the words that you can identify in this list", the other 2 tests had me find synonyms of said word and calculated my score based on how many synonyms I correctly identified, and the last test had me identify words by giving me the option to select "yes/no" if I knew the word, and if I selected "yes" it would randomly ask me for the synonym of said word and calculated my score based on if I answered correctly or not (a combo of the two styles that was supposed to keep you honest).

Overall, these tests were a decent way to figure out how many words you know for someone who is brainrotted and has a short attention span, but not perfect. I do have some problems with the method of choice. One major quip I had with these tests are the words they chose are ones that I would find in a Jane Austen novel (bring back "thistledown" as a common word to be used in our vernacular again). I know that these words exist in the English language, and the question "how many words do you know in the English language" is a very broad one, and choosing words at random in a dictionary would also yield weird Austen-esque words as well. Obviously these methods aren't foolproof. It is almost impossible to test for the thousands of words that are contained in the English language without taking up 8 hours of your day or your entire lifetime for that matter. So, while flawed (maybe) the question of "how many" was answered!

\section{So, does this even matter?}

While having a big vocabulary is impressive and will earn you brownie points in a fifth grade classroom, I fail to see how testing to see "how many words I know in English" has anything to do with me grasping entire concepts or building a foundation of knowledge about the world. Knowing individual words doesn't really enhance my ability to experience and perceive the world in a great manner. While it is nice, being literate $\ne$ being knowledgeable about the world. Recall that to have knowledge, one must make successful contact with reality, and through this contact we can understand, interpret, and interact with the world. Whether or not this statement is entirely truthful (which is me blissfully turning a blind eye to the counter-arguments one could pose), knowing a handful of words can't by itself successfully build my repertoire of experiences. Rather being able to perceive, experience, and make judgments, decisions, and living as a free agent (if we are truly "free") gives me fulfillment and teaches me about the world and gives me some semblance of having knowledge. I wouldn't consider myself a particularly "smart" individual (says the graduate student in computer science), yet I do believe that I have had a life full of many different experiences that do not have any correlation to the amount of words I know.

% ==================
% DOCUMENT ENDS HERE
% ==================
\end{document}
